
\bta{2007B}

\begin{enumerate}
	%\renewcommand{\labelenumi}{\arabic{enumi}.}
	% A(\Alph) a(\alph) I(\Roman) i(\roman) 1(\arabic)
	%设定全局标号series=example	%引用全局变量resume=example
	%[topsep=-0.3em,parsep=-0.3em,itemsep=-0.3em,partopsep=-0.3em]
	%可使用leftmargin调整列表环境左边的空白长度 [leftmargin=0em]
	\item
（$30'$）考虑一维阶梯势$V(x)=\begin{cases}V_0,&x>0(V_0>0)\\ 0,&x<0\end{cases}$，设粒子从右边向左边入射，试求反射系数和透射系数。


\banswer{
	从右侧存在传播入射波要求总能量 $E>V_0$。令
	\[
		q=\frac{\sqrt{2m(E-V_0)}}{\hbar}\quad(x>0),
		\qquad
		k=\frac{\sqrt{2mE}}{\hbar}\quad(x<0).
	\]
	取
	$\psi_{>}=e^{-iqx}+r e^{iqx}$、$\psi_{<}=t e^{-ikx}$。在 $x=0$ 匹配 $\psi$ 与 $\psi'$，得到
	\[
		r=\frac{q-k}{q+k},\qquad t=\frac{2q}{q+k}.
	\]
	按概率流之比，
	\[
		\boxed{R=\left(\frac{q-k}{q+k}\right)^2,\qquad
		T=\frac{k}{q}|t|^2=\frac{4qk}{(q+k)^2}},
		\qquad R+T=1.
	\]
}


\item 
（$30'$）电子处于沿$+z$方向、大小为$B$的均匀磁场中。设$t=0$时刻电子自旋沿$+y$方向。
\begin{enumerate}
	%\renewcommand{\labelenumi}{\arabic{enumi}.}
	% A(\Alph) a(\alph) I(\Roman) i(\roman) 1(\arabic)
	%设定全局标号series=example	%引用全局变量resume=example
	%[topsep=-0.3em,parsep=-0.3em,itemsep=-0.3em,partopsep=-0.3em]
	%可使用leftmargin调整列表环境左边的空白长度 [leftmargin=0em]
	\item
试求$t=0$时电子自旋波函数；

\item 
试分别求$t>0$时，电子自旋沿$+x,+y,+z$方向的概率。
	
\end{enumerate}

\banswer{
	在 $S_z$ 表象中，$t=0$ 时沿 $+y$ 的自旋态为
	\[
		\boxed{\chi(0)=|+y\rangle=\frac1{\sqrt2}
		\begin{pmatrix}1\\i\end{pmatrix}}.
	\]
	把 Zeeman 哈密顿量写成 $H=(\hbar\Omega/2)\sigma_z$；$\Omega$ 的符号包含电子磁矩的符号。于是
	\[
		\chi(t)=\frac1{\sqrt2}
		\begin{pmatrix}e^{-i\Omega t/2}\\i e^{i\Omega t/2}\end{pmatrix}.
	\]
	分别投影到 $|+x\rangle$、$|+y\rangle$、$|+z\rangle$，得
	\[
		\boxed{P_{+x}(t)=\frac{1-\sin\Omega t}{2},\qquad
		P_{+y}(t)=\frac{1+\cos\Omega t}{2},\qquad
		P_{+z}(t)=\frac12}.
	\]
	若采用相反号的 Zeeman 哈密顿量，只需令 $\Omega\to-\Omega$；物理上即进动方向反转。
}


\item 
（$30'$）
粒子在势场 $\quad V(x)=\left\{\begin{array}{cl}V_{0} \delta(x), & |x|<a \\ \infty, & |x| \geq a\end{array} \quad\right.$ 中运动 $\left(V_{0}>0\right)$ 。 试求系统能级和能级方程。


\banswer{
	势阱关于原点对称，故本征态可分为奇、偶宇称。
	\begin{itemize}
		\item 奇态在 $x=0$ 处为零，不受 $\delta$ 势影响：
		\[
			\psi_n^{(-)}(x)\propto\sin\frac{n\pi x}{a},\qquad
			k_n=\frac{n\pi}{a},\qquad
			E_n^{(-)}=\frac{\hbar^2n^2\pi^2}{2ma^2},\quad n=1,2,\ldots.
		\]
		\item 偶态在 $0<x<a$ 可写为 $\psi(x)=C\sin[k(a-x)]$，并作偶延拓。$\delta$ 势导致导数跃变
		\[
			\psi'(0^+)-\psi'(0^-)=\frac{2mV_0}{\hbar^2}\psi(0).
		\]
		代入偶态后得到能级方程
		\[
			\boxed{k\cot(ka)=-\frac{mV_0}{\hbar^2}},
			\qquad
			\boxed{E^{(+)}=\frac{\hbar^2k^2}{2m}}.
		\]
	\end{itemize}
	这两组能级合起来构成系统的全部能谱；$V_0>0$ 只把偶宇称能级向上推移。
}


\newpage
\item 
（$30'$）设系统哈密顿算符为$\hat{H}=\frac{\hat{p}^2}{2m}+V(\vec{r})$，粒子处于归一化的束缚定态$\psi_n$中。试证明Virial定理：
$$\langle \psi_n|\frac{\hat{p}^2}{2m}|\psi_n\rangle=\frac{1}{2}\langle\psi_n |\vec{r}\cdot\nabla V(\vec{r})|\psi_n\rangle$$

\banswer{
	定义厄米的位力算符
	\[
		G=\frac12(\vec r\cdot\vec p+\vec p\cdot\vec r).
	\]
	利用正则对易关系可得
	\[
		[H,G]=i\hbar\left[-2\frac{p^2}{2m}+\vec r\cdot\nabla V(\vec r)\right].
	\]
	在归一化束缚定态中，$\langle G\rangle$ 不随时间改变；由 Heisenberg 方程，
	\[
		0=\frac{d\langle G\rangle}{dt}
		=\frac{i}{\hbar}\langle[H,G]\rangle.
	\]
	因此
	\[
		2\left\langle\frac{p^2}{2m}\right\rangle
		=\langle\vec r\cdot\nabla V\rangle,
	\]
	即
	\[
		\boxed{\left\langle\frac{p^2}{2m}\right\rangle
		=\frac12\langle\vec r\cdot\nabla V(\vec r)\rangle}.
	\]
}


\item 
（$30'$）一维谐振子系统哈密顿量为$\hat{H}_0=\frac{\hat{p}^2}{2m}+\frac{1}{2}m\omega^2x^2$，设受到微扰$H'=-\lambda \hat{p}_x^4$的作用。试求对等$n$个谐振子能级的一级微扰修正。\\
（已知矩阵元$\langle n'|\hat{x}|n\rangle=\sqrt{\frac{\hbar}{2m\omega}}(\sqrt{n+1}\delta_{n',n+1}+\sqrt{n}\delta_{n',n-1})$）

\banswer{
	动量可写为
	\[
		p=i\sqrt{\frac{m\hbar\omega}{2}}(a^\dagger-a).
	\]
	用升降算符展开，或利用谐振子坐标与动量分布的对称性，可得
	\[
		\langle n|p^4|n\rangle
		=\left(\frac{m\hbar\omega}{2}\right)^2
		(6n^2+6n+3)
		=\frac34m^2\hbar^2\omega^2(2n^2+2n+1).
	\]
	故一级能量修正为
	\[
		\boxed{E_n^{(1)}=\langle n|H'|n\rangle
		=-\frac{3\lambda}{4}m^2\hbar^2\omega^2
		(2n^2+2n+1)}.
	\]
}


\end{enumerate}
